\documentclass[10pt,letterpaper,final]{report}
\usepackage[margin=.75in]{geometry}
\usepackage[utf8]{inputenc}
\usepackage{amsmath}
\usepackage{amsfonts}
\usepackage{amssymb}
\usepackage{amsthm}
\renewcommand\qedsymbol{$\blacksquare$}
\usepackage{enumerate}
\usepackage{enumitem}
\usepackage{hyperref}
\usepackage{pdfpages}
\usepackage{graphics}
\usepackage{graphicx}
\usepackage{tikz}
\usepackage{tikz-qtree}
\usepackage{forest}
\usepackage{float}
\usetikzlibrary{automata,arrows}
\usetikzlibrary{shapes.multipart, positioning}
\usepackage{mdframed}
\usepackage[
  backend=biber,
  style=numeric-comp,
  sorting=none
]{biblatex}
\addbibresource{references.bib}

\usepackage{minted}
\usemintedstyle{algol_nu}
\usepackage{xltabular}
\usepackage{pdflscape}
\usepackage{tcolorbox}
\usepackage{array}
\usepackage{todonotes}
\setuptodonotes{inline}  % render \todo notes as full-width blocks in the text flow, not margin boxes

\newcommand{\reason}[1]{&& \text{#1}}

% Based on template from COSC-417 at Towson University, originally authored by Marius Zimand

% To input a script, adapt the following:
% \inputminted[breaklines, breakanywhere, fontsize=\footnotesize, frame=single, fontsize=\footnotesize]{python}{script.py}

\begin{document}
\begin{center}
  \fbox{
    \vbox{
      \begin{flushleft}
        Jonathan Llovet \\  % authors' names
        EN.625.252.81 \\  %class
        2026-10-18 \\  % date
        Module 7 - Homework 7 \\ % ID
      \end{flushleft}
      \center{\Large{\textbf{Linear Algebra - Determinants}}}
      %\end{mdframed}
    } % end vbox
  } % end fbox
  \vline
\end{center}

\medskip
\noindent\textbf{Problem 1:}
\medskip

Compute the determinants of the following matrices. \emph{Note:} See Appendix B if needed for a refresher on complex numbers.

\begin{enumerate}[label=(\alph*)]
  \item $\begin{bmatrix}
            8 & 0  \\
            3 & -1
          \end{bmatrix}$
  \item $\begin{bmatrix}
            2i & 3  \\
            4  & 6i
          \end{bmatrix}$, where $i = \sqrt{-1}$
  \item $\begin{bmatrix}
            i  & 2 + i & 0   \\
            -1 & 3     & 2i  \\
            0  & -1    & 1-i
          \end{bmatrix}$, where $i = \sqrt{-1}$
\end{enumerate}

\medskip
\noindent\textbf{Answer:}

\begin{enumerate}[label=(\alph*)]
  \item The definition of the determinant of a $2 \times 2$ matrix $A = \begin{bmatrix}
            a & b \\
            c & d
          \end{bmatrix}$ is $\det A = ad - bc$.

        Here then we can calculate: \begin{align*}
          \det \begin{bmatrix}
                 8 & 0  \\
                 3 & -1
               \end{bmatrix} & = (8)(-1) - (0)(3) \\[1ex]
                                 & = -8 - 0       \\
                                 & = -8
        \end{align*}
  \item This case is the same as the one above, except that we must perform some multiplication of complex numbers. \begin{align*}
          \det \begin{bmatrix}
                 2i & 3  \\
                 4  & 6i
               \end{bmatrix} & = (2i)(6i) - (3)(4) \\[1ex]
                                 & = 12i^2 - 12    \\[1ex]
                                 & = 12(-1) - 12   \\[1ex]
                                 & = -12 - 12      \\[1ex]
                                 & = -24           \\[1ex]
        \end{align*}
  \item In this third case, we have to take advantage of cofactor expansion. (I have added braces to make the algebra easier to trace visually.)
        \begin{align*}
          \det A & = \det \begin{bmatrix}
                            i  & 2 + i & 0   \\
                            -1 & 3     & 2i  \\
                            0  & -1    & 1-i
                          \end{bmatrix}  \reason{Given}                                                                                                                                                        \\[1ex]
                 & = \underbrace{a_{11}C_{11}} + \underbrace{a_{12}C_{12}} + \underbrace{a_{13}C_{13}} \reason{By definition of determinant \cite[Ch. 3.1, p.173]{LayEtAl2020LinearAlgebraAndItsApplications}} \\[1ex]
                 & = \underbrace{iC_{11}} + \underbrace{(2 + i)C_{12}} + \underbrace{(0)C_{13}} \reason{By substitution}                                                                                       \\[1ex]
                 & = \underbrace{i(-1)^{1+1}\det A_{11}} + \underbrace{(2 + i)(-1)^{1+2}\det A_{12}} \reason{By substitution and simplification of mult. by 0}                                                 \\[1ex]
                 & = \underbrace{i(-1)^{2}\det A_{11}} + \underbrace{(2 + i)(-1)^{3}\det A_{12}} \reason{By simplification}                                                                                    \\[1ex]
                 & = \underbrace{i\det A_{11}} + \underbrace{(2 + i)(-1)\det A_{12}} \reason{By simplification}                                                                                                \\[1ex]
                 & = \underbrace{i\det \begin{bmatrix}3 & 2i \\-1 & 1-i\end{bmatrix}} + \underbrace{(2 + i)(-1)\det \begin{bmatrix}-1 & 2i \\0 & 1-i\end{bmatrix}} \reason{By simplification}                  \\[1ex]
                 & = \underbrace{i(3(1-i) - (-1(2i)))} + \underbrace{(2 + i)(-1)(-1(1-i)-2i(0))} \reason{By evaluation of the determinants}                                                                    \\[1ex]
                 & = \underbrace{i(3 - 3i + 2i)} + \underbrace{(2 + i)(-1)(-1 + i)} \reason{By simplification}                                                                                                 \\[1ex]
                 & = \underbrace{i(3 - i)} + \underbrace{(2 + i)(1 - i)} \reason{By simplication}                                                                                                              \\[1ex]
                 & = \underbrace{3i - i^2} + \underbrace{2 - 2i + i - i^2} \reason{By multiplication}                                                                                                          \\[1ex]
                 & = \underbrace{3i - (-1)} + \underbrace{2 - i - (-1)} \reason{By evaluation of $i^2$}                                                                                                        \\[1ex]
                 & = \underbrace{3i + 1} + \underbrace{2 - i + 1} \reason{By multiplication by $-1$}                                                                                                           \\[1ex]
                 & = 2i + 4 \reason{By addition}                                                                                                                                                               \\[1ex]
                 & = 2(2 + i) \reason{By factoring, commutativity}
        \end{align*}
\end{enumerate}


\pagebreak
\noindent\textbf{Problem 2:}
\medskip

(Continuation of the prior assignment.) Let the matrices $A = \begin{bmatrix}
    3  & -2 & -2 \\
    -1 & 1  & 1  \\
    3  & -1 & -2
  \end{bmatrix}$ and $B = \begin{bmatrix}
    1  & a & 0  \\
    -1 & b & 1  \\
    2  & c & -1
  \end{bmatrix}$ be inverses of each other. Find the value of $b$ using Cramer's Rule. Compare your answer to your prior assignment.

\medskip
\noindent\textbf{Answer:}

\pagebreak
\noindent\textbf{Problem 3:}
\medskip

Let $T: R^3 \to R^3$ be the linear transformation determined by the matrix $A = \begin{bmatrix}
    a & 0 & 0 \\
    0 & b & 0 \\
    0 & 0 & c
  \end{bmatrix}$, where $a$, $b$, and $c$ are positive real numbers. Let $B$ be the unit sphere defined by $x^{2}_{1} + x^{2}_{2} + x^{2}_{3} = 1$

\begin{enumerate}[label=(\alph*)]
  \item Show that the image $T(B)$ is the ellipsoid with the equation $\frac{x^{2}_{1}}{a^2} + \frac{x^{2}_{2}}{b^2} + \frac{x^{2}_{3}}{c^2} = 1$
  \item Use determinants and the fact that the volume of the unit sphere is $\frac{4}{3}\pi$ to determine the volume formula for the ellipsoid in part (a).
\end{enumerate}

\medskip
\noindent\textbf{Answer:}

\pagebreak

\section*{Source Code}
The full source code, including LaTeX, tooling, other artifacts, and git history is available on my Github repository for the class here:
\begin{center}
  \url{https://github.com/jllovet/jhu-en-625-252-fa26-linear-algebra}
\end{center}

\printbibliography

\end{document}
